\SourceLeaf{97}valorem per $\pi\pi l$ dividere oporteret, denotante $l$ longitudinem penduli simplicis per unitatem temporis vibrantis, ita ut generaliter habeatur $TM=\frac{K}{ttl\cos i}$, vel pro casu nostro $TM=\frac{K}{ttl}$.

\Article{8}
Si experimenta huius generis in acubus magneticis instituuntur ad filum verticale suspensis, reactio, quam torsio exserit, in experimentis subtilioribus haud negligenda erit. Distinguamus in tali filo duos diametros horizontales, alterum $D$ in termino inferiori, ubi acus adnexa est, axi magnetico acus parallelum, alterum $E$ in termino superiori, ubi filum fixum est, ipsi $D$ parallelum in statu detorsionis. Supponamus, $E$ facere cum meridiano magnetico angulum $v$, contra axem magneticum vel $D$ angulum $u$, eritque experientia duce vis torsionis, proxime saltem, angulo $v-u$ proportionalis: statuemus itaque momentum, quo haec vis angulum $u$ ipsi $v$ aequalem reddere nititur. $=(v-u)\theta$. Iam quum momentum vis magneticae terrestris ad angulum $u$ minuendum sit $=TM\sin u$, conditio aequilibrii continetur in aequatione $(v-u)\theta=TM\sin u$, quae eo plures solutiones reales admittet, quo minor est $\theta$ respectu ipsius $TM$: quatenus autem hic tantummodo de valoribus parvis ipsius $u$ agitur, tuto eius loco hanc adoptare licet $(v-u)\theta=TMu$ sive $\frac vu=\frac{TM}{\theta}+1$. In apparatibus nostris terminus fili superior brachio horizontali mobili adnexus est, quod portat indicem in peripheria circuli in gradus divisi incedentem. Etiamsi itaque error collimationis (i. e. divisio cui respondet valor $v=0$) nondum satis exacte cognitus sit, tamen iste index differentiam binorum valorum ipsius $v$ monstrat: perinde alia apparatus pars differentiam inter valores ipsius $u$ statui aequilibrii respondentes summa praecisione subministrat, patetque, valorem ipsius $\frac{TM}{\theta}+1$ e divisione differentiae inter duos valores ipsius $v$ per differentiam inter valores respondentes ipsius $u$ obtineri Quatenus inter experimenta ad hunc finem instituenda temporis intervallum aliquanto longius praeterlabitur, necesse erit, si summa praecisio desideratur, ut variationis diurnae declinationis magneticae ratio habeatur, quod facile fit adiumento observationum simultanearum in secundo apparatu, in quo fili terminus superior intactus conservatur: vix opus est monere, distantiam inter ambos apparatus tantam esse debere, ut sensibiliter se mutuo turbare nequeant.

\expandafter\def\csname HistoricalNoteText9\endcsname{Proprie haec fila partialia non sunt vere simplicia, sed tantummodo talia, qualia a mercatoribus non neta venduntur.}
\SourceLeaf{98}
\expandafter\def\csname HistoricalNoteText8\endcsname{Utraeque divisiones a laeva versus dextram crescunt.}
Ut quantam subtilitatem huiusmodi observationes admittant eluceat, exemplum e diario adscribimus. Observatae sunt $1832$ Sept $22$, salvis erroribus collimationis, declinationes $u$ atque anguli $v$ sequentes\HistoricalNote{8}:
\begin{center}
\begin{tabular}{c|c||c|c||c}
&&\multicolumn{2}{c||}{Acus prima}&Acus secunda\\
Exp.&tempus&$u$&$v$&$u$\\\hline
I&$9^h33'$ matut.&$+0^0\ 4'19''5$&$300^0$&$+0^0\ 2'12''1$\\
II&$9\ 57$&$-0\ 0\ 19{,}6$&$240$&$+0\ 1\ 37{,}7$\\
III&$10\ 16$&$-0\ 4\ 40{,}5$&$180$&$+0\ 1\ 18{,}8$
\end{tabular}
\end{center}
Sunt itaque declinationes acus primae ad statum primae observationis reductae hae
\begin{center}
\begin{tabular}{crr}
I.&$u=+0^0\ 4'19''5$&$v=300^0$\\
II&$+0\ 0\ 14{,}8$&$240$\\
III&$-0\ 3\ 47{,}2$&$180$
\end{tabular}
\end{center}
Hinc prodit valor fractionis $\frac{TM}{\theta}$ e combinatione observationum
\begin{center}
\begin{tabular}{cr}
I et II $.\ .\ .\ .\ .$&$881{,}7$\\
II et III $.\ .\ .\ .\ .$&$891{,}5$\\
I et III $.\ .\ .\ .\ .$&$886{,}6$
\end{tabular}
\end{center}

Variationes declinationis magneticae diurnae per torsionem in ratione unitatis ad $\frac{n}{n+1}$ minuuntur, statuendo $\frac{TM}{\theta}=n$, quae mutatio, si filis tam parvae torsionis, qualem exemplum praecedens exhibet, utimur, pro insensibili haberi potest. Quod vero attinet ad tempus oscillationum (infinite parvarum), e principiis dynamicis facile concluditur, hoc in ratione unitatis ad $\BareRoot{\frac{n}{n+1}}$ per torsionem minui. Proprie haec referuntur ad casum eum, ubi $v=0$: formulae vero generaliter valerent, si statueremus $\frac{TM\cos u^0}{\theta}=n$, denotando per $u^0$ valorem ipsius $u$ aequilibrio respondentem: sed differentia prorsus insensibilis erit.

\Article{9}
Coëfficiens $\theta$ principaliter pendet a longitudine, crassitie et materia fili; insuper in filis metallicis aliquantulum a temperatura, in bombycinis a statu hygrometrico: contra in illis (forsanque etiam in his, \textit{dum sunt simplicia}) haud qua\SourceLeaf{99}quam a pondere, quo onerantur, pendere videtur. Aliter vero se habet res in filis bombycis compositis, quales ad acus graviores ferendas adhibere oportet: in his $\theta$ cum pondere appenso augetur, multo tamen minor manet valore ipsius $\theta$ pro filo metallico eiusdem longitudinis eidemque ponderi ferendo apto. Ita e. g. per methodum prorsus similem ei, quam in art. praec. tradidimus (sed in alio filo aliaque acu), inventus est valor ipsius $n=597{,}4$, dum filum portabat acum cum sola supellectile ordinaria, ubi pondus integrum erat $496{,}2$ grammatum; contra $=424{,}8$, quum pondus usque ad $710{,}8$ grammata auctum esset, sive erat in casu primo $\theta=0{,}0016740\,TM$, in casu secundo $\theta=0{,}0023542\,TM$. Filum, cuius longitudo est $800$ millimetrorum, compositum est e $32$ filis simplicibus\HistoricalNote{9}, quae singula $30$ fere grammata tuto portant, atque ita ordinata sunt, ut aequalem tensionem patiantur. Ceterum verisimile est, valorem ipsius $\theta$ constare e parte constante et parte ponderi proportionali, atque partem constantem aequalem fieri aggregato valorum ipsius $\theta$ pro singulis filis simplicibus. In hac hypothesi (per experimenta hactenus nondum satis confirmata) pars constans pro exemplo allato invenitur $=0{,}0001012\,TM$, adeoque valor ipsius $\theta$ filo simplici respondens $=0{,}00000316\,TM$. Adiumento valoris ipsius $TM$ mox eruendi ex hac hypothesi colligitur, reactionem fili simplicis per arcum radio aequalem $(57^0\ 18')$ torsi aequivalere gravitati milligrammatis in vectem longitudinis circiter $\frac1{17}$ millimetri prementis.

\Article{10}
Si corpus oscillans est acus simplex figurae regularis massaeque homogeneae, momentum inertiae $K$ per methodos notas calculari potest. E. g. si corpus est parallelepipedum rectangulum, cuius latera sunt $a,b,c$, densitas $=d$, et proin massa $q=abcd$, momentum inertiae respectu axis per centrum transeuntis laterique $c$ paralleli erit $=\frac1{12}(aa+bb)q$: et quum in acubus magneticis talis formae latus, cui axis magneticus parallelus est, $a$, longe maior esse soleat latitudine $b$, pro experimentis crassioribus adeo sufficiet, statuere $K=\frac1{12}aaq$. At in experimentis subtilioribus, etiam ubi acus simplex adhibetur, suppositionem gratuitam massae perfecte homogeneae formaeque perfecte regularis aegre admit\SourceLeaf{100}teremus, et pro experimentis nostris, ubi non acus simplex, sed acus cum supellectile complicatiore iuncta oscillat, rem per talem calculum expedire omnino impossibile est, adeoque de alio modo, momentum $K$ maxima praecisione determinandi, cogitare oportuit.

Cum acu coniungebatur virga lignea transversalis, a qua pendebant duo pondera aequalia, per cuspides acutissimas in puncta virgae $A,B$ prementia: haec puncta erant in recta horizontali, in eodem plano verticali cum axe suspensionis, et utrimque inde aeque distantia. Denotando massam utriusque ponderis per $p$, distantiam $AB$ per $2r$, per accessionem huius apparatus momentum $K$ augebitur quantitate $C+2prr$, ubi $C$ est aggregatum momenti inertiae virgae respectu lineae suspensionis atque momentorum ponderum respectu axium verticalium per cuspides et centra gravitatis transeuntium. Si itaque oscillationes tum acus non oneratae, tum acus in duabus distantiis diversis oneratae, puta pro $r=r'$ atque $r=r''$ observatae, temporaque oscillationum (ad infinite parvas reductarum et ab effectu torsionis purgatarum) resp. $=t,t',t''$ inventa sunt, e combinatione aequationum
\[
\begin{aligned}
TMtt&=\pi\pi K\\
TMt't'&=\pi\pi(K+C+2pr'r')\\
TMt''t''&=\pi\pi(K+C+2pr''r'')
\end{aligned}
\]
tres incognitae $TM,K$ et $C$ erui poterunt. Praecisionem adhuc maiorem assequemur, si observatis oscillationibus pro pluribus valoribus ipsius $r$, puta pro $r=r',r'',r'''$ etc. respondentibus temporibus $t',t'',t'''$ etc., per methodum quadratorum minimorum duas incognitas $x,y$ ita determinamus, ut satisfiat quam proxime aequationibus
\[
\begin{aligned}
t'&=\BareRoot{\frac{r'r'+y}{x}}\\
t''&=\BareRoot{\frac{r''r''+y}{x}}\\
t'''&=\BareRoot{\frac{r'''r'''+y}{x}}\ \text{etc.}
\end{aligned}
\]
quo facto habebimus
\[
\begin{aligned}
TM&=2\pi\pi px\\
K+C&=2py
\end{aligned}
\]

\SourceLeaf{101}
Circa hanc methodum sequentia adhuc monere convenit.

I. Quoties acus non nimis laevigata adhibetur, sufficit, virgam ligneam simpliciter illi imponere. Quoties autem superficies acus perlaevis est, ut frictio impedire nequeat, quominus virga super illa gliscere possit, necesse est, quo totus apparatus ad instar unius corporis rigidi moveatur, virgam apparatui reliquo firmius adstringere. In utroque vero casu prospiciendum est, ut puncta $A,B$ sint satis exacte in recta horizontali.

II. Quum complexus talium experimentorum aliquot horas postulet, variabilitas intensitatis magnetismi terrestris intra hoc temporis spatium, siquidem summa praecisio desideratur, haud negligenda est. Quocirca antequam eliminatio suscipiatur, tempora observata ad valorem constantem ipsius $T$, e. g. ad valorem medium experimento primo respondentem, reducere oportet. Ad hunc finem observationibus simultaneis in alia acu (perinde ut in art.\,8.) opus est, quae si tempus unius oscillationis pro temporibus mediis singulorum experimentorum resp. prodiderunt $=u,u',u'',u'''$ etc., ad calculum loco valorum observatorum $t',t'',t'''$ etc. resp. adhibendi sunt hi
\[
\frac{ut'}{u'},\qquad\frac{ut''}{u''},\qquad\frac{ut'''}{u'''}\quad\text{etc.}
\]

III. Simile monitum valet circa variabilitatem ipsius $M$, a variatione temperaturae, si quae inter experimenta locum habuit, oriundam. Sed patet, reductionem modo adscriptam iam per se hanc correctionem implicare, si utraque acus aequali temperaturae mutationi subiecta fuerit, et perinde a tali mutatione afficiatur.

IV. Quoties tantummodo de valore ipsius $TM$ eruendo agitur, manifesto experimentum primum superfluum est. Attamen utile erit, experimentis acu onerata factis statim adiungere aliud acu non onerata, ut simul valor ipsius $K$ prodeat, qui experimentis alio tempore eadem acu instituendis substrui possit, quum manifesto hic valor invariatus maneat, etiamsi $T$ et $M$ lapsu temporis mutationem subire possint.

\Article{11}
Ad maiorem illustrationem huius methodi e magna copia applicationum exemplum unum hic adscribimus. Ecce conspectum numerorum, quos experimenta $1832$ Sept. $11$ instituta prodiderunt.
